\section{DeepSeek V4.1 and architecture level prefill reuse}
\label{sec:C03:frontier}
The September 2026 DeepSeek-V4.1-Flash report makes prefill reuse an architectural choice. Its Section 2.2 derives upper-layer global KV projections from a middle-layer hidden state, while retaining layerwise sliding-window computation. Section 2.3 distinguishes full, reindex, and reuse modes for compressed sparse attention; the report also describes a four-bit main KV cache~\cite{rewrite_deepseek41}. These mechanisms differ from merely allocating an ordinary dense cache more efficiently. Their correctness and quality depend on a model trained with those structural choices, so they cannot be applied as a transparent replacement for every earlier checkpoint.

In a schematic causal encoder-decoder with $L$ layers, $n$ prompt tokens, and bounded replay of $w$ tokens in the upper half, a linear-in-length layer-cost model gives a prefill ratio of $(nL/2+wL/2)/(nL)=(1+w/n)/2$. This is an analytical model of the stated decomposition, not a measured speedup. If $n$ is much larger than $w$, the ratio approaches one half. If $n=w$, the modeled reduction disappears. Actual projection, attention, index construction, and scheduling costs can move the result. The formula nevertheless identifies why a long-document cache miss and a short tool-response turn are not interchangeable workloads.

The same case qualifies a blanket invalidation rule while requiring a distinction between structural reuse and approximation. The report's Section 3.2.2 explicitly truncates sliding-window replay to the recent window and accepts approximate states; exact reconstruction would require a longer dependency history. Cached global projections and bounded sliding-window replay therefore have different equivalence claims~\cite{rewrite_deepseek41}. In a conventional dense decoder, later representations may depend on the entire prefix. A trained sharing architecture changes that computation graph, but does not make every reconstruction exact. Bind each reuse invariant, approximation, and quality test to the checkpoint and implementation. A cache produced for a different projection path, compression mode, or precision is not compatible merely because its token positions match.

A discriminating experiment compares cold prefill, exact eligible-prefix reuse, and repeated short tool turns while keeping the checkpoint, quality tasks, and external tool behavior fixed. Inspect physical cache bytes and actual recomputed tokens, not only the advertised context length. Include edits to an early permission constraint and questions requiring precise old identifiers. A model-specific compression method can preserve aggregate long-context scores while failing these rare distinctions. Conversely, a specialized architecture may achieve substantial resource savings without sacrificing the tested distinctions. The result must remain attached to the actual mechanism and workload rather than a generic claim of superior memory.




